\documentclass{ltjsarticle}
\usepackage{amsmath}
\usepackage{geometry}
\geometry{body={165mm,230mm}}

\def\version{3.4}
\begin{document}

\begin{center}
  {\LARGE 楽しい運動計測実習: 基礎編(ドラム演奏バージョン)}
\end{center}
\begin{flushright}
\today
\end{flushright}

\section*{実習での注意}
\begin{itemize}
\item 計測実験の時には、「データをとっては確認」を繰り返すのが王道。
  計測データを全て計測してから計測時の不備に気づいて、全データの取り直しにならないように注意すること。
\item 以下の実験では, 原則として全員のデータを取って解析してください。個人差が無いかを確認することは重要です。
\item 実験の操作, 準備は役割分担をして行うこと。時間内に実験が終わるように常に意識する。
\item 解析プログラムは各人で作成後, お互いにその出力が同じになっているか相互に確認すること。バグ取りは重要です。
\item 取得データはデータ格納専用の外部HDD等に専用フォルダを作って格納すること。計測用パソコンのローカルディスクには放置しないでください。身元不明ファイルになってしまい, 後で困ります。
\item 取得データのファイル名は「日付\_グループ名と被験者番号\_性別(M or F)\_実験番号と条件.csv」とすること(Aグループ1人目の実験1の条件(a)時の例: 250318\_A01\_M\_1a.csv)。
% \item 運動計測実験では, 身長・体重等の身体パラメータがわかったほうが考
%   察に役立つこともありますが, 体重等を秘密にしたい方は無理に測らなくて
%   もOKです。
\end{itemize}

\section{実験1: 負荷と筋電位}

\subsection{必要な機器等}

筋電計測用パソコン, 筋電センサ一式, 実験で手に持つ重り(2kg, 4kg, 6kg), 被験者チェックシート(筋電, 反射マーカ, セッションシート), ビデオカメラ, メジャー, 筋電センサ貼付用テープ

\subsection{実験}

\begin{enumerate}
  \item 身長と体重を測定する。
  \item 手首関節, 肘関節, 肩関節に反射マーカを貼付しなさい(実験2で計測)。またマーカ間距離もチェックシートに記録すること。
  \item 上腕二頭筋, 上腕三頭筋, 尺側手根屈筋, 尺側手根伸筋に筋電センサを貼付しなさい(尺側手根屈筋と尺側手根伸筋は実験2で計測)。チェックシートに確認チェックを入れること。
  \item 最大随意収縮の筋電位(MVC)を測定しなさい。
  \item 筋電位の周波数分布はおおむね 5 Hz から 500 Hz である。したがって, 計測のサンプリング周波数は 1000 Hz 以上である必要がある。今回は1,926 Hzで計測する。
  \item 上腕は鉛直下向き, 前腕を前に水平に出した状態で, いろいろな重り(重りなし, 2kg, 4kg, 6kg)を 10 秒間持ったときの上肢の筋電位(EMG)を計測しなさい。計測データはネーミングポリシーに従って保存すること。
\end{enumerate}  

実験終了後は筋電計と反射マーカを必ず取り外すこと。

また取得データの処理については別紙「楽しい運動計測実習, データ解析編」を見ること。

\subsection{プレゼンテーション/レポート}

負荷と, 上腕二頭筋, 上腕三頭筋の筋電位の大きさとの関係を解析し, その結果及び自分なりの発見をレポート(and/or プレゼンテーション)にまとめなさい。

計測データに対しては, いろいろなデータ処理や可視化の方法があり得る。データ背後にある法則性についての仮説を考え, それを検証するためにはどのようなデータ処理をしてどのようなグラフを作れば良いかよく考えること。以下は主な考察ポイント。

\begin{enumerate}
  \item 平均的な筋活動はおもりの重量とともに線形に増加するか, 非線形に増加するか
  \item 重りを持って姿勢維持を続けた場合, 筋活動には時間に伴う変化はあるか
  \item 以上については, 平均値だけでなく標準偏差についても議論すること
  \item (optional) 負荷に伴う周波数分布の変化はあるだろうか
  \item その他発見はあるか
\end{enumerate}

プレゼンテーション等の資料を作る時には, 作ったグラフをやみくもに並べるのではなく, 自分の主張を根拠とともに明確に伝えるには, どのグラフ(何についての関係性)を, どのように(縦軸や横軸のレンジ等)作れば良いかを良く考えること。


\section{実験2: ドラム演奏の動作計測}

\subsection{必要な機器等}

メトロノーム(パソコンかスマホ利用), 筋電センサ一式, 筋電計測用パソコン, モーションキャプチャ一式, モーションキャプチャ用パソコン, 同期ユニット, 被験者チェックシート(筋電, 反射マーカ, セッションシート), メジャー, 両面テープ, サージカルテープ

\subsection{実験}

\begin{enumerate}
\item 身長と体重を測定する(学会発表等では重要となるが今回は省略する)。
\item 手首関節, 肘関節, 肩関節に反射マーカを貼付しなさい。またマーカ間距離も記録すること。
\item 上腕二頭筋, 上腕三頭筋, 尺側手根屈筋, 尺側手根伸筋に筋電センサを貼付しなさい。
\item 最大随意収縮の筋電位(MVC)を測定しなさい。
\item 筋電位の周波数分布はおおむね 5 Hz から 500 Hz である。
したがって, 計測のサンプリング周波数は 1000 Hz 以上である必要がある。今回は1,926 Hzで計測する。
%\item 各関節位置の軌跡はOptiTrackで抽出する。
%\item 得られたデータは自分のPCに転送して解析すること。
\item ドラム演奏を行う際の上肢関節軌道と筋電位(EMG)を計測しなさい。
      打叩時には, はじめの姿勢で肘を下げ, 肘の屈曲で腕が矢状面内で動くような姿勢を取り, 矢状面内の運動になるように気をつけること。
      また手の甲を上に向けた状態でスティックを握り, 指は動かさずに打叩すること。
\item 実験条件は以下のとおりとする。
  \begin{enumerate}
    \item 中程度の強さ(mf)で, 120bpm(回/分)で約20秒間台を叩く % テンポと強弱の基準
    \item 非常に強く(ff), 120bpm(回/分)で約20秒間台を叩く % 強く叩いた場合
    \item 非常に弱く(pp), 120bpm(回/分)で約20秒間台を叩く % 弱く叩いた場合
    \item 中程度の強さ(mf)で, 240bpm(回/分)で約20秒間台を叩く % 速く叩いた場合
    \item 中程度の強さ(mf)で, 120bpm(回/分)で約20秒間台を叩く
  \end{enumerate}
  (a)〜(e)の順番で計測を行いなさい。ただし, 条件(b), (c)は被験者により順序を変えることで順序効果の影響を小さくする。
\item 各条件のテンポに合わせたメトロノームが鳴っている間に, タスク毎で30秒 ~ 1分程度の練習を行い, 慣れてきた時点で計測を開始しなさい。
\item セッションシートに計測条件や気付き(トラブルや被験者の運動の特徴)をメモすること。また, 明らかなストロークミスやメトロノームとのズレがあった場合は計測をし直すこと。
\item 計測データはネーミングポリシーに従って保存すること。
\end{enumerate}
実験終了後は筋電計と反射マーカを必ず取り外すこと。

\subsection{プレゼンテーション/レポート}

以下をレポート(and/or プレゼンテーション)にまとめなさい。
\begin{enumerate}
  \item 各筋活動は, ドラム演奏の打叩運動に対してどのタイミングでおきるだろうか(力学の問題として捉えて考えること。)
  \item ドラム演奏の強弱を変えたとき, 活動する筋肉や活動の大きさには変化があるだろうか。その理由はなにか。
  \item ドラム演奏のテンポを変えるとき, 筋活動は周期が変わるのみか。各筋活動のタイミング(位相)に変化はあるだろうか。一周期あたりの活動時間に変化はあるだろうか。その理由はなにか。
  \item その他, 自分なりの考察や発見を説明しなさい。
\end{enumerate}

\end{document}